\section{FISH}
\label{sec:method}
FISH takes a ROS~2 application as it is launched and returns, for one run of it, the model of Fig.~\ref{fig:model}. It has four components, used in this order: (i) \emph{instrumentation}, tracepoints added to the ROS~2 client libraries, and the CUDA profiler; (ii) \emph{acquisition}, an entry point and a daemon that run the unmodified application under the instrumentation and record one complete trace; (iii) \emph{extraction}, which turns the trace into the model; and (iv) \emph{FISH-eye}, a viewer for the model (Fig.~\ref{fig:screens}). FISH currently targets ROS~2 Humble.

\noindent\textbf{Instrumentation.} \texttt{ros2\_tracing}~\cite{bedard-2022-ros2tracing} records each node, interface and callback as it is created and, for timer, subscription and service callbacks, each execution; the thread comes from the tracer's context. To its 14 events FISH adds 28 that close the gaps of Section~\ref{sec:goal}: start and end events for client and waitable callbacks, the intra-process subscription among them; introspection events for each executor's type and thread count and each callback's group; link events that tie each publication and request to the callback that issued it; finalize events, so that a reused handle is not mistaken for its predecessor; and the rclpy counterparts of all of these. GPU work is recorded by CUPTI through \texttt{nsys}, aligned to the LTTng clock at ingest and bound through correlation ids to the API call and, through the thread's active callback, to the callback that issued it. The OS scheduler's context switches are recorded as well and split each execution into on-core, preempted and blocked time.

\noindent\textbf{Acquisition.} The instrumented libraries replace the stock ones; the application is not touched. The FISH entry point wraps \texttt{ros2 launch}, reads the launch description and starts every container and node that will use the GPU under \texttt{nsys} from the outset, while a daemon polls \texttt{/proc} and restarts under the profiler any GPU process started outside the launch. The daemon then waits until the set of loaded components and started processes has been stable over several polls, marks the system stable, and starts the input: a rosbag replay, or for a live system timed commands issued from the system-stable instant.

\noindent\textbf{Extraction.} Initialisation events yield the vertices of Fig.~\ref{fig:model}a and the message and service edges the wiring declares; execution events find the callback that actually issued each, add the state and join edges from the trace patterns of Fig.~\ref{fig:model}c and attach the measurements; every edge is lifted along the containment (Fig.~\ref{fig:model}b), and each GPU-submitting callback receives the DAG, shapes and conditional graph of Fig.~\ref{fig:model}d. Phases are delimited from the trace itself: initialisation ends when every periodic timer has fired once, warm-up when every recurring callback has completed its first execution, and what follows is the steady state. Finally, the model is checked against what the traced binaries create, e.g., a node constructor must yield one node, seven entities and seven callbacks.

\noindent\textbf{FISH-eye} renders the model at any layer of Fig.~\ref{fig:model}a, at mixed granularity through the expansion operator, and grouped by namespace (Fig.~\ref{fig:screens}).
